\chapter{页错误}
\label{CH:PGFAULTS}
%%

当访问页表中无映射的虚拟地址，或映射的 \lstinline{PTE_V} 标志已清除的地址，或映射的权限位（\lstinline{PTE_R}、\lstinline{PTE_W}、\lstinline{PTE_X}、\lstinline{PTE_U}）禁止当前操作时，RISC-V CPU 会触发页错误异常。RISC-V 区分三种页错误：加载页错误（由加载指令引起）、存储页错误（由存储指令引起）和指令页错误（由取指执行指令引起）。\lstinline{scause} 寄存器指示页错误的类型，\indexcode{stval} 寄存器包含无法翻译的地址。

页表与页错误的组合是一个强大的工具。页表在内核虚拟地址与物理地址之间提供了一层间接层，使内核能够控制地址空间的结构和内容。页错误则允许内核拦截加载和存储操作，并通过修改页表来实时指定这些引用所指向的数据。内核可以利用这些能力提高效率：例如，写时复制fork允许内核在父进程和子进程之间透明地共享内存，避免复制双方都不会写入的页面。应用程序员也能从中受益。一种可能性是内存映射文件，内核通过分页机制将文件内容映射到应用程序的地址空间中，在响应页错误时透明地读取文件块。另一种是惰性内存分配，允许程序申请巨大的虚拟地址空间，但只需为实际读写的页面支付分配物理内存的成本。xv6 仅将页错误用于一个目的：惰性分配。

在继续之前，请阅读 {\tt kernel/sysproc.c} 中的函数 {\tt sys\_sbrk()} 和 {\tt kernel/vm.c} 中的 {\tt vmfault}。在 {\tt kernel/trap.c} 和 {\tt kernel/vm.c} 中搜索对 {\tt vmfault} 的调用。

\section{惰性分配}
\label{sec:lazy}

xv6 的 \indextext{lazy allocation} 包含两部分。首先，当应用程序通过调用带有 \lstinline{SBRK_LAZY} 标志的 \lstinline{sbrk} 请求内存时，内核仅记录大小的增加，不为新范围的虚拟地址分配物理内存或创建 PTE。其次，当在这些新地址之一上发生页错误时，内核分配一页物理内存并将其映射到页表中。内核透明地实现惰性分配：应用程序无需修改即可受益。

惰性分配对应用程序很方便，因为它们不必准确预测需要多少内存。例如，应用程序可能要处理输入，但事先不知道输入有多大。使用惰性分配，应用程序可以为最坏情况申请内存，但不必为此付费：内核无需为应用程序从未使用的页面做任何工作。

此外，如果应用程序要求大幅增加地址空间，那么没有惰性分配的 \lstinline{sbrk} 会很昂贵：如果应用程序申请 1 GB 内存，内核必须分配并清零 262,144 个 4096 字节的物理页面。惰性分配允许将这种成本分散到时间上。但另一方面，惰性分配会产生页错误的额外开销，这涉及用户/内核态切换。操作系统可以通过每次页错误分配一批连续页面而非单个页面，以及为这种页错误专门优化内核进入/退出代码来降低这种成本（尽管 xv6 两者都不做）。

另一方面，当为惰性分配页面发生页错误时，内核可能发现没有空闲内存可分配。在这种情况下，内核无法向应用程序返回内存不足错误，只能杀死应用程序。对于希望在分配失败时收到错误的应用程序，xv6 允许通过使用 \lstinline{SBRK_EAGER} 标志调用 \lstinline{sbrk} 来立即分配内存。

\section{代码}

系统调用 {\tt sbrk(n)} 通过 {\tt n} 字节增加（如果 {\tt n} 为负则减少）进程的内存大小，并返回新分配区域的起始地址（即旧大小）。内核实现是 \lstinline{sys_sbrk}
\lineref{kernel/sysproc.c:/^sys_sbrk/}。

如果应用程序指定 \lstinline{SBRK_EAGER}，系统调用由函数
\lstinline{growproc}
\lineref{kernel/proc.c:/^growproc/}
实现。\lstinline{growproc} 调用 \lstinline{uvmalloc}。
\lstinline{uvmalloc}
\lineref{kernel/vm.c:/^uvmalloc/}
使用 {\tt kalloc} 分配物理内存，将分配的内存清零，并使用 {\tt mappages} 将 PTE 添加到用户页表。

如果应用程序惰性分配内存，\lstinline{sys_sbrk}
只将进程的大小
(\lstinline{myproc()->sz}) 增加 {\tt n} 并返回旧大小；它不分配物理内存或向进程页表添加 PTE。

当进程对缺少有效页表映射的虚拟地址进行加载或存储时，CPU 将引发 \indextext{page-fault exception}。
\lstinline{usertrap} 检查此情况
\lineref{kernel/trap.c:/page fault/}
并调用 \lstinline{vmfault}
\lineref{kernel/vm.c:/^vmfault/}
来处理页错误。\lstinline{vmfault}
检查出错地址是否在 {\tt sbrk} 先前授予的区域内，
使用 {\tt kalloc} 分配一页物理内存，
将分配的页面清零，并使用
{\tt mappages} 将 PTE 添加到用户页表。Xv6 在 PTE 中设置新页面的 \lstinline{PTE_W}、\lstinline{PTE_R}、\lstinline{PTE_U} 和 \lstinline{PTE_V} 标志。然后，\lstinline{usertrap} 在导致错误的指令处恢复进程。因为 PTE 现在有效，重新执行的加载或存储指令将不会出错。

如果应用程序使用 {\tt sbrk} 释放内存，\lstinline{sys_sbrk}
调用 \lstinline{shrinkproc}，后者调用 \lstinline{uvmdealloc}。实际工作由 {\tt uvmunmap} \lineref{kernel/vm.c:/^uvmunmap/} 完成，它使用 {\tt walk} 查找 PTE。由于某些页面可能从未被进程使用，因此从未由 {\tt vmfault} 分配，{\tt uvmunmap} 跳过没有 \lstinline{PTE_V} 标志的 PTE。如果 PTE 映射有效，{\tt uvmunmap} 调用 {\tt kfree} 释放它引用的物理内存。

请注意，xv6 使用进程页表不仅告诉硬件如何映射用户虚拟地址，还作为分配给该进程的物理内存页面的唯一记录。这就是在 {\tt uvmunmap} 中释放用户内存时需要检查用户页表的原因。

\section{现实世界：写时复制（COW）fork}

许多内核（尽管不是 xv6）使用页错误来实现
\indextext{copy-on-write (COW) fork}。{\tt fork} 系统调用承诺子进程看到的内存的初始内容与 fork 时父进程的内存相同。实现这一点的一种方法是将父进程的整个内存复制到为子进程新分配的物理内存中；这就是 xv6 所做的。复制可能很慢，如果子进程可以共享父进程的物理内存，效率会更高。然而，这种简单实现不会奏效，因为它会导致父进程和子进程通过写入共享的栈和堆来相互干扰执行。

写时复制fork通过适当使用页表权限和页错误使父进程和子进程安全地共享物理内存。基本计划是让父进程和子进程最初共享所有物理页面，但每个都将其映射为只读（\lstinline{PTE_W} 标志清除）。然后父进程和子进程都可以从共享物理内存中读取。如果任何一方写入共享页面，RISC-V CPU 会触发页错误异常。支持 COW 的内核会响应分配新物理内存页面并将共享页面复制到新页面中。然后内核会更改错误进程的页表中的相关 PTE 以指向副本并允许读写，然后在导致错误的指令处恢复错误进程。因为 PTE 现在允许写入，重新执行的存储指令将不会出错，并将修改页面的私有副本而不是共享页面。

写时复制需要簿记来帮助决定何时可以释放物理页面，因为每个页面可以根据 fork、页错误、exec 和 exit 的历史被不同数量的页表引用。这种记账允许一个重要的优化：如果进程产生存储页错误且物理页面仅从该进程的页表引用，则不需要复制。

写时复制使 \lstinline{fork} 更快，因为 \lstinline{fork} 不需要复制内存。某些内存稍后写入时必须复制，但通常大部分内存永远不需要复制。一个常见例子是
\lstinline{fork} 后跟 \lstinline{exec}：
\lstinline{fork} 之后可能写入几页，
但随后子进程的 \lstinline{exec} 释放从父进程继承的大部分内存。写时复制 \lstinline{fork} 消除了复制这些内存的需要。
此外，COW fork 是透明的：
应用程序无需修改即可受益。

\section{现实世界：按需分页}

另一个利用页错误的广泛使用的功能是
\indextext{demand paging}。在 \lstinline{exec} 系统调用中，xv6 在启动应用程序之前将所有应用程序的文本和数据加载到内存中。由于应用程序可能很大，从磁盘读取需要时间，这种启动成本对用户来说可能很明显。为了减少启动时间，现代内核最初不将可执行文件加载到内存中，而只是创建用户页表并将所有 PTE 标记为无效。内核启动程序运行；每次程序第一次使用页面时，发生页错误，内核响应从磁盘读取页面内容并将其映射到用户地址空间。像 COW fork 和惰性分配一样，内核可以透明地为应用程序实现此功能。

在计算机上运行的程序可能需要比计算机拥有的 RAM 更多的内存。为了优雅地应对，操作系统可以实现 \indextext{paging to disk}。想法是只在 RAM 中存储一小部分用户页面，其余存储在磁盘上的 \indextext{paging area} 中。内核将对应于存储在分页区域中（因此不在 RAM 中）的内存的 PTE 标记为无效。如果应用程序尝试使用已 \emph{分页出去} 到磁盘的页面之一，应用程序将产生页错误，页面必须被 \emph{分页进来}：内核陷阱处理程序将分配一页物理 RAM，将页面从磁盘读入 RAM，并修改相关 PTE 以指向 RAM。

如果需要分页进入页面，但没有空闲物理 RAM 怎么办？在这种情况下，内核必须首先通过将其分页出去或 \emph{驱逐} 到磁盘上的分页区域，并将引用该物理页面的 PTE 标记为无效来释放物理页面。驱逐是昂贵的，因此如果它不频繁，分页性能最好：如果应用程序只使用其内存页面的一个子集，且这些子集的并集适合 RAM。这个属性通常被称为具有良好的引用局部性。与许多虚拟内存技术一样，内核通常以实现为对应用程序透明的方式实现分页到磁盘。

计算机通常在很少有或没有 \emph{空闲} 物理内存的情况下运行，无论硬件提供多少 RAM。例如，云提供商将许多客户多路复用到单台机器上以有效使用其硬件。另一个例子是用户在少量物理内存的智能手机上运行许多应用程序。在这种情况下，分配页面可能需要首先驱逐现有页面。因此，当空闲物理内存稀缺时，分配是昂贵的。

当空闲内存稀缺且程序只积极使用其分配内存的一部分时，惰性分配和按需分页特别有利。这些技术还可以避免分配或加载页面但从未使用或在可以使用之前被驱逐时浪费的工作。

\section{现实世界：内存映射文件}

结合分页和页错误异常的其他功能包括自动扩展栈和 \indextext{memory-mapped files}，这是程序使用 \texttt{mmap} 系统调用映射到其地址空间的文件，以便程序可以使用加载和存储指令读写它们。

% "virtual" memory, eviction, page-in
% lazy allocation
% auto stack expansion
% guard pages
% mmap files
% cow fork
% shared text, shared libraries
% demand paging of text
% virtual machine migration
% distributed shared memory
% fast IPC
% zero-copy write
% DPDK
% (unified block / page cache)

%%
\section{练习}
%%

\begin{enumerate}

\item 编写一个用户程序，通过调用
\lstinline{sbrk(1)}
将其地址空间增加一个字节。
运行该程序并调查程序调用
\lstinline{sbrk}
之前和调用
\lstinline{sbrk}
之后的页表。
内核分配了多少空间？新内存的
PTE
包含什么？

\item 实现 COW fork。

\item 实现 {\tt mmap}。

\end{enumerate}
